\section{Participant Demographics}
\label{sec:eval-demographics}

Before reporting the user-experience and document-quality findings,
this section profiles the participant sample. It describes the
matched within-subjects sample of $n = 69$ valid participants
($p_{0001} \geq 1$) on which all subsequent findings rest. Four
screened records were excluded from all later analyses for
incomplete or invalid participation, yielding the analysis sample
reported in Chapter~\ref{ch:methodology}. The profile is presented
descriptively here; its implications for the generalisability of the
findings are reserved for the Discussion chapter.

\subsection{Numeric and Ordinal Variables}

Table~\ref{tab:demographics-numeric} reports the descriptive
statistics for the four numeric or ordinal variables collected in the
demographics block: total session duration in seconds, and three
five-point Likert self-reports of prior expertise (GDPR knowledge,
ROPA familiarity, chatbot familiarity).

\begin{table}[htbp]
\centering
\caption{Descriptive statistics for numeric and ordinal demographic
variables ($n = 69$). Duration is in seconds and aggregates the full
study session across all three modes; the three expertise variables
are five-point Likert self-reports (1 = no familiarity, 5 = expert).
The duration minimum of $-1$ is a sentinel value for one
non-responding record on this single variable; the median is a more
robust point estimate.}
\label{tab:demographics-numeric}
\small
\begin{tabular}{@{}l r r r r r r r r@{}}
\toprule
Variable & $n$ & Mean & SD & Min & Q1 & Median & Q3 & Max \\
\midrule
Duration (s)         & 69 & 1865.7 & 1354.2 & $-1$ & 840 & 2013 & 2708 & 5788 \\
GDPR knowledge       & 69 & 2.81  & 1.15  & 1    & 2   & 3    & 4    & 5    \\
ROPA familiarity     & 69 & 1.45  & 0.76  & 1    & 1   & 1    & 2    & 4    \\
Chatbot familiarity  & 69 & 4.36  & 0.73  & 2    & 4   & 4    & 5    & 5    \\
\bottomrule
\end{tabular}
\end{table}

The median session duration is $2013$~s ($\approx 33.6$~min) per
participant across the full session, consistent with the protocol
target of $\approx 90$~min total once spread across three modes and
the rating instruments. ROPA familiarity is the lowest of the four
variables (mean $1.45$, median $1$), confirming that the sample
qualifies as ROPA-novice in line with the inclusion criteria. Chatbot
familiarity is correspondingly high (mean $4.36$, median $4$). GDPR
knowledge sits near the scale midpoint (mean $2.81$).

\subsection{Categorical Variables}

Table~\ref{tab:demographics-categorical} reports counts and shares
for the five categorical demographic variables: gender, age (binned
by year), highest completed degree, current professional status, and
field of study or work. Three participants did not report gender; the
share columns therefore use the per-variable response base.

\begin{table}[htbp]
\centering
\caption{Categorical demographic variables ($n = 69$ unless noted).
Shares are computed within variable.}
\label{tab:demographics-categorical}
\small
\begin{tabular}{@{}l l r r@{}}
\toprule
Variable & Level & $n$ & Share \\
\midrule
\multirow{3}{*}{Gender ($n = 66$ reporters)}
  & Male           & 41 & 62.1\,\% \\
  & Female         & 24 & 36.4\,\% \\
  & Other          &  1 & \phantom{0}1.5\,\% \\
\midrule
\multirow{4}{*}{Age (years)}
  & 24             & 18 & 26.1\,\% \\
  & 25             & 18 & 26.1\,\% \\
  & 26             & 10 & 14.5\,\% \\
  & Other (21--40) & 23 & 33.3\,\% \\
\midrule
\multirow{3}{*}{Degree}
  & Bachelor       & 60 & 88.2\,\% \\
  & Master         &  6 & \phantom{0}8.8\,\% \\
  & Other          &  2 & \phantom{0}2.9\,\% \\
\midrule
\multirow{4}{*}{Profession}
  & Student        & 56 & 81.2\,\% \\
  & Employed       & 10 & 14.5\,\% \\
  & PhD            &  2 & \phantom{0}2.9\,\% \\
  & Staatsexamen   &  1 & \phantom{0}1.5\,\% \\
\midrule
\multirow{5}{*}{Field}
  & IT             & 41 & 59.4\,\% \\
  & Business (BWL) & 16 & 23.2\,\% \\
  & Other          &  6 & \phantom{0}8.7\,\% \\
  & Research       &  4 & \phantom{0}5.8\,\% \\
  & HR             &  2 & \phantom{0}2.9\,\% \\
\bottomrule
\end{tabular}
\end{table}

The sample is dominated by Bachelor-level students of IT or Computer
Science. Of the $66$ participants who reported gender, $62.1\,\%$
identified as male. The full distribution of demographics and
expertise variables is visualised in Figure~\ref{fig:demographics}.
This profile — IT-leaning, chatbot-fluent, ROPA-novice — frames the
results reported in \S\ref{sec:eval-userstudy}--\S\ref{sec:eval-judge}
and is taken up again in the Discussion chapter when generalisability
is addressed.

\begin{figure}[htbp]
\centering
\includegraphics[width=\linewidth]{demographics.png}
\caption{Demographic and expertise distributions: histograms of age
and session duration, count plots for gender, degree, profession, and
field, and a grouped bar chart of the three expertise self-reports
(GDPR, ROPA, chatbot) on the five-point Likert scale.}
\label{fig:demographics}
\end{figure}
